\section{Conclusion}\label{sec:concl}
This paper proposed a blockchain-enabled Zero Trust architecture for identity management in Industrial IoT, built around the observation that a plant needs multi-party revocation authority, verification proportional to consequence, and an enforcement path that does not wait for consensus. Three smart contracts anchor identity, tiered policy and quorum-governed revocation; a fast path decides locally against replicated state and anchors Merkle roots of its decisions, while resources whose misuse is costly are decided on-chain with an anomaly gate and a strict mode; and batch revocation supports firmware recalls. We proved immediate revocation on the audited path, a $\rho\Delta$ exposure bound on the fast path, replay resistance, attestation unforgeability, key-reuse prevention and governance safety and liveness, and showed exactly when a spoofing attacker is accepted.

On a synthetic industrial fleet, the audited design accepted no request after revocation where token and CRL baselines accepted 33--147 on average; the gate and strict mode reduced full-spoof acceptance at the two most critical tiers to 0\% and stealthy-mimicry acceptance to 60.9\% and 13.6\%, at a cost of 2.2--5.2\% denied legitimate requests; batch revocation saved about 40\% of gas, with about 35k gas per device; anchoring cost 54 gas per decision at 1{,}000 decisions per root; and quorum revocation for 5--21 voters cost 313k--800k gas, with no coalition below quorum succeeding in 1{,}200 attempts. With the behavioural signal taken from real TON\_IoT traffic instead, the full policy accepted under 1\% of seen-type attacker requests at the critical tiers and 6.6--11.9\% of held-out-type requests on average (worst type 33.9--41.6\% at T1--T2), while an adaptive attacker that sends only its best-scoring flows was accepted 29--77\% of the time; telemetry alone was a weak signal for most devices. On a second dataset (WUSTL-IIOT-2018, SCADA) a supervised detector gave the same picture, whereas an unsupervised one cut attack acceptance at the two critical tiers to 4--8\% only at the price of denying 12--18\% of legitimate requests there. A centralised online policy decision point matches the zero exposure of the audited path, so the case for the design rests on shared revocation authority and an externally anchored audit trail, not on exposure alone. On a four-validator Besu QBFT network, audited calls were final in about one 2\,s block, a revocation was visible on a second validator in 1.9\,s and denied all 30 later audited requests, a three-of-five quorum revocation took 4--8\,s, sustained throughput was only 11--18 decisions per second on one two-vCPU host, and the chain halted safely when two of four validators were stopped. The remaining timing numbers are compute-bound, the posture and network signals remain synthetic, and an attacker that keeps its behaviour within normal variation is accepted at every tier. Whether the approach holds up in a real plant is a question for deployment across several hosts with real protocols, devices and telemetry, which we regard as the next step together with threshold-signed oracles and privacy protection for exposed scores. Code, result files and generation scripts are released at \href{https://github.com/irfanrazanaqvi/BCZT_IDM_IoT_experments_code}{github.com/irfanrazanaqvi/BCZT\_IDM\_IoT\_experments\_code}; the TON\_IoT data must be obtained from its owners.

\balance
\begin{thebibliography}{99}\fontsize{7.5}{8.8}\selectfont
\setlength{\itemsep}{0pt}
\setlength{\parskip}{0pt}
\fontsize{7.5}{8.8}\selectfont
\setlength{\itemsep}{0pt}
\setlength{\parskip}{0pt}
\bibitem{stouffer2023sp80082}
K. Stouffer \emph{et al.}, ``Guide to Operational Technology (OT) Security,'' NIST Special Publication 800-82 Rev.~3, National Institute of Standards and Technology, Sep. 2023. \href{https://doi.org/10.6028/NIST.SP.800-82r3}{[Online]}

\bibitem{nist800207}
S. Rose, O. Borchert, S. Mitchell, and S. Connelly, ``Zero Trust Architecture,'' NIST Special Publication 800-207, National Institute of Standards and Technology, 2020. \href{https://doi.org/10.6028/NIST.SP.800-207}{[Online]}

\bibitem{fernandez2024critical}
E. B. Fernandez and A. Brazhuk, ``A critical analysis of zero trust architecture (ZTA),'' \emph{Computer Standards \& Interfaces}, vol. 89, Art. no. 103832, 2024. \href{https://doi.org/10.1016/j.csi.2024.103832}{[Online]}

\bibitem{phiayura2023comprehensive}
P. Phiayura and S. Teerakanok, ``A comprehensive framework for migrating to zero trust architecture,'' \emph{IEEE Access}, vol. 11, pp. 19487--19511, 2023. \href{https://doi.org/10.1109/ACCESS.2023.3248622}{[Online]}

\bibitem{song2025ztot}
M. W. Song, T. Nguyen, and C. Irvine, ``A zero trust architecture for critical operational technology systems,'' in \emph{Proc. 58th Hawaii Int. Conf. Syst. Sci. (HICSS)}, 2025. \href{https://doi.org/10.24251/HICSS.2025.849}{[Online]}

\bibitem{zanasi2024flexible}
C. Zanasi, S. Russo, and M. Colajanni, ``Flexible zero trust architecture for the cybersecurity of industrial IoT infrastructures,'' \emph{Ad Hoc Networks}, vol. 156, Art. no. 103414, 2024. \href{https://doi.org/10.1016/j.adhoc.2024.103414}{[Online]}

\bibitem{w3cdid2022}
M. Sporny \emph{et al.}, ``Decentralized Identifiers (DIDs) v1.0,'' W3C Recommendation, Jul. 2022. \href{https://www.w3.org/TR/did-core/}{[Online]}

\bibitem{hosseini2023blockchain}
S. M. Hosseini, J. Ferreira, and P. C. Bartolomeu, ``Blockchain-based decentralized identification in IoT: an overview of existing frameworks and their limitations,'' \emph{Electronics}, vol. 12, no. 6, p. 1283, 2023. \href{https://doi.org/10.3390/electronics12061283}{[Online]}

\bibitem{aguru2025ssi}
A. D. Aguru and S. B. Erukala, ``A blockchain-based self-sovereign identity management framework for Internet of Things,'' in \emph{Recent Trends in Intelligent Systems and Next Generation Wireless Communication (IIWCS 2024)}, Lecture Notes in Networks and Systems, vol. 1329. Singapore: Springer, 2025, pp. 299--309. \href{https://doi.org/10.1007/978-981-96-4741-5_26}{[Online]}

\bibitem{folino2025puf}
G. Folino, A. Forestiero, and G. Papuzzo, ``Self-sovereign identification of IoT devices by using physically unclonable functions and blockchain,'' in \emph{Numerical Computations: Theory and Algorithms (NUMTA 2023)}, Lecture Notes in Computer Science, vol. 14478. Cham: Springer, 2025. \href{https://doi.org/10.1007/978-3-031-81247-7_26}{[Online]}

\bibitem{zhai2025chebyshev}
Z. Zhai \emph{et al.}, ``A lightweight authentication method for Industrial Internet of Things based on blockchain and Chebyshev chaotic maps,'' \emph{Future Internet}, vol. 17, no. 8, Art. no. 338, 2025. \href{https://doi.org/10.3390/fi17080338}{[Online]}

\bibitem{essaid2025iiotbc}
M. Essaid and H. Ju, ``Blockchain solutions for enhancing security and privacy in Industrial IoT,'' \emph{Applied Sciences}, vol. 15, no. 12, Art. no. 6835, 2025. \href{https://doi.org/10.3390/app15126835}{[Online]}

\bibitem{liu2024dissecting}
C. Liu, R. Tan, Y. Wu, Y. Feng, Z. Jin, F. Zhang, Y. Liu, and Q. Liu, ``Dissecting zero trust: research landscape and its implementation in IoT,'' \emph{Cybersecurity}, vol. 7, no. 1, Art. no. 20, 2024. \href{https://doi.org/10.1186/s42400-024-00212-0}{[Online]}

\bibitem{nie2025zeroTrust}
S. Nie, J. Ren, R. Wu, P. Han, Z. Han, and W. Wan, ``Zero-trust access control mechanism based on blockchain and inner-product encryption in the Internet of Things in a 6G environment,'' \emph{Sensors}, vol. 25, no. 2, p. 550, 2025. \href{https://doi.org/10.3390/s25020550}{[Online]}

\bibitem{salim2026ztblociot}
M. M. Salim, M. Kim, S. K. Singh, and J. H. Park, ``Zero-trust blockchain-enabled framework for scalable and secure IoT networks (ZT-BlocIoT),'' \emph{Future Generation Computer Systems}, vol. 175, Art. no. 108093, 2026. \href{https://doi.org/10.1016/j.future.2025.108093}{[Online]}

\bibitem{mukta2026ztdelegation}
R. Mukta, S. Pal, K. Chowdhury, M. Hitchens, H. Paik, and S. S. Kanhere, ``Zero trust-driven access control delegation using blockchain,'' \emph{Blockchain: Research and Applications}, vol. 7, no. 1, Art. no. 100319, 2026. \href{https://doi.org/10.1016/j.bcra.2025.100319}{[Online]}

\bibitem{jamwal2026survey}
E. Jamwal, P. Srivastava, and P. S. Chakraborty, ``Blockchain-enabled zero-trust framework for secure cloud data sharing: survey,'' in \emph{Proc. 2026 Int. Conf. Intell. Syst. Eng., Secured Syst. Cybersecurity (ICISESSC)}, 2026. \href{https://doi.org/10.1109/ICISESSC68634.2026.11542695}{[Online]}

\bibitem{farhan2026ztfsf}
M. Farhan \emph{et al.}, ``ZTF-SF: a zero-trust blockchain-enabled IIoT framework for secure smart factories,'' in \emph{Proc. 2026 2nd Int. Conf. Comput. Intell. Approaches Appl. (ICCIAA)}, 2026. \href{https://doi.org/10.1109/ICCIAA68481.2026.11543806}{[Online]}

\bibitem{potturi2026hybrid}
K. Potturi, B. Gautam, C. S. Pareek, A. Gupta, S. C. Chava, and B. S. Katta, ``A hybrid AI and blockchain-enabled zero-trust security framework for cloud, edge, and IoT systems,'' in \emph{Proc. 2026 Int. Conf. Comput. Theory Wireless Commun. (ICCTWC)}, 2026. \href{https://doi.org/10.1109/ICCTWC68241.2026.11582964}{[Online]}

\bibitem{wang2026medicaliot}
Z. Wang, Y. Guo, and Y. Guo, ``Blockchain-based zero-trust continuous authentication scheme for medical IoT,'' \emph{Journal of Parallel and Distributed Computing}, vol. 217, Art. no. 105333, 2026. \href{https://doi.org/10.1016/j.jpdc.2026.105333}{[Online]}

\bibitem{james2024authz}
M. James, T. Newe, D. O'Shea, and G. D. O'Mahony, ``Authentication and authorization in zero trust IoT: A survey,'' in \emph{Proc. Irish Signals and Systems Conf. (ISSC)}, 2024. \href{https://doi.org/10.1109/ISSC61953.2024.10603175}{[Online]}

\bibitem{kumar2024industry40}
R. Kumar and R. Sharma, ``A comprehensive approach to industry 4.0 security: blockchain and dynamic access control integration,'' in \emph{Proc. 2024 6th Int. Conf. Computational Intelligence and Communication Technologies (CCICT)}, 2024, pp. 477--484. \href{https://doi.org/10.1109/CCICT62777.2024.00080}{[Online]}

\bibitem{azad2024verify}
M. A. Azad, S. Abdullah, J. Arshad, H. Lallie, and Y. H. Ahmed, ``Verify and trust: a multidimensional survey of zero-trust security in the age of IoT,'' \emph{Internet of Things}, vol. 27, Art. no. 101227, 2024. \href{https://doi.org/10.1016/j.iot.2024.101227}{[Online]}

\bibitem{nahar2024survey}
N. Nahar, K. Andersson, O. Schel\'en, and S. Saguna, ``A survey on zero trust architecture: applications and challenges of 6G networks,'' \emph{IEEE Access}, vol. 12, pp. 94753--94764, 2024. \href{https://doi.org/10.1109/ACCESS.2024.3425350}{[Online]}

\bibitem{ren2025aiot}
J. Ren, J. Zhang, Y. Ren, and J. Xu, ``Blockchain-based self-sovereign identity management mechanism in AIoT environments,'' \emph{Electronics}, vol. 14, no. 19, p. 3954, 2025. \href{https://doi.org/10.3390/electronics14193954}{[Online]}

\bibitem{philipp2023daxiot}
A. Philipp and A. K\"upper, ``DAXiot: A decentralized authentication and authorization scheme for dynamic IoT networks,'' arXiv:2307.06919, 2023. \href{https://doi.org/10.48550/arXiv.2307.06919}{[Online]}

\bibitem{mazzocca2024evoke}
C. Mazzocca, A. Acar, S. Uluagac, and R. Montanari, ``EVOKE: Efficient revocation of verifiable credentials in IoT networks,'' in \emph{Proc. 33rd USENIX Security Symp.}, Philadelphia, PA, USA, Aug. 2024, pp. 1279--1295. \href{https://www.usenix.org/system/files/usenixsecurity24-mazzocca.pdf}{[Online]}

\bibitem{kumar2026privacy}
A. Kumar, B. Sharma, R. K. Kaushal, and I. Ben Dhaou, ``Privacy-preserving identity management system leveraging blockchain technology,'' in \emph{Proc. 15th IEEE Int. Conf. Commun. Syst. Netw. Technol. (CSNT)}, 2026, p. 625 ff. \href{https://doi.org/10.1109/CSNT69054.2026.11502058}{[Online]}

\bibitem{usama2026multilayer}
M. Usama, A. Aziz, N. Alasbali, N. Alturki, M. Hanif, and M. Ur Rehman, ``Blockchain-enabled identity management for IoT: a multi-layered defense against adversarial AI,'' \emph{Scientific Reports}, vol. 16, Art. no. 4371, 2026. \href{https://doi.org/10.1038/s41598-026-35208-y}{[Online]}

\bibitem{tom2025flcti}
A. K. Tom, A. Khraisat, T. Jan, M. Whaiduzzaman, T. D. Nguyen, and A. Alazab, ``Survey of federated learning for cyber threat intelligence in Industrial IoT: Techniques, applications and deployment models,'' \emph{Future Internet}, vol. 17, no. 9, Art. no. 409, 2025. \href{https://doi.org/10.3390/fi17090409}{[Online]}

\bibitem{shujaa2025enhancing}
W. Shujaa, M. Alanzi, and S. Sankaranarayanan, ``Enhancing IoT security through blockchain integration,'' \emph{Frontiers in Computer Science}, vol. 7, Art. no. 1670473, 2025. \href{https://doi.org/10.3389/fcomp.2025.1670473}{[Online]}

\bibitem{hassan2023oracle}
A. Hassan, I. Makhdoom, W. Iqbal, A. Ahmad, and A. Raza, ``From trust to truth: Advancements in mitigating the blockchain oracle problem,'' \emph{J. Netw. Comput. Appl.}, vol. 217, Art. no. 103672, 2023. \href{https://doi.org/10.1016/j.jnca.2023.103672}{[Online]}

\bibitem{wei2024quality}
Z. Wei \emph{et al.}, ``Survey on quality assurance of smart contracts,'' \emph{ACM Comput. Surv.}, vol. 57, no. 2, Art. no. 32, 2024. \href{https://doi.org/10.1145/3695864}{[Online]}

\bibitem{ta2024gascost}
M. T. Ta and T. Q. Do, ``A study on gas cost of Ethereum smart contracts and performance of blockchain on simulation tool,'' \emph{Peer-to-Peer Networking and Applications}, vol. 17, no. 1, pp. 200--212, 2024. \href{https://doi.org/10.1007/s12083-023-01598-3}{[Online]}

\bibitem{hasan2023smartcontract}
M. R. Hasan, A. Alazab, S. B. Joy, M. N. Uddin, M. A. Uddin, A. Khraisat, I. Gondal, W. F. Urmi, and M. A. Talukder, ``Smart contract-based access control framework for Internet of Things devices,'' \emph{Computers}, vol. 12, no. 11, Art. no. 240, 2023. \href{https://doi.org/10.3390/computers12110240}{[Online]}

\bibitem{haque2024scalable}
E. U. Haque, A. Shah, J. Iqbal, S. S. Ullah, R. Alroobaea, and S. Hussain, ``A scalable blockchain based framework for efficient IoT data management using lightweight consensus,'' \emph{Scientific Reports}, vol. 14, Art. no. 7841, 2024. \href{https://doi.org/10.1038/s41598-024-58578-7}{[Online]}

\bibitem{peepliwal2024ztab}
A. K. Peepliwal, H. M. Pandey, S. Prakash, A. A. Mahajan, S. S. Chowhan, V. Kumar, and R. Sharma, ``A prototype model of zero trust architecture blockchain with EigenTrust-based practical Byzantine fault tolerance protocol to manage decentralized clinical trials,'' \emph{Digital Communications and Networks}, 2024. \href{https://www.sciencedirect.com/science/article/pii/S2096720924000459}{[Online]}

\bibitem{madhwal2025gasless}
Y. Madhwal \emph{et al.}, ``Empowering autonomous IoT devices in blockchain through gasless transactions,'' \emph{Blockchain: Research and Applications}, vol. 6, no. 2, Art. no. 100257, 2025. \href{https://doi.org/10.1016/j.bcra.2024.100257}{[Online]}

\bibitem{ren2026ditdp}
L. Ren, H. Li, Y. Liu, D. Liao, M. Zhang, and H. Jin, ``Identity trust management based on differential privacy for internet of vehicles,'' \emph{Scientific Reports}, vol. 16, Art. no. 20921, 2026. \href{https://doi.org/10.1038/s41598-026-50273-z}{[Online]}

\bibitem{alsaedi2020ton}
A. Alsaedi, N. Moustafa, Z. Tari, A. Mahmood, and A. Anwar, ``TON\_IoT telemetry dataset: A new generation dataset of IoT and IIoT for data-driven intrusion detection systems,'' \emph{IEEE Access}, vol. 8, pp. 165130--165150, 2020. \href{https://doi.org/10.1109/ACCESS.2020.3022862}{[Online]}

\bibitem{moustafa2021network}
N. Moustafa, ``A new distributed architecture for evaluating AI-based security systems at the edge: Network TON\_IoT datasets,'' \emph{Sustainable Cities and Society}, vol. 72, Art. no. 102994, 2021. \href{https://doi.org/10.1016/j.scs.2021.102994}{[Online]}

\bibitem{booij2022ton}
T. M. Booij, I. Chiscop, E. Meeuwissen, N. Moustafa, and F. T. H. den Hartog, ``ToN\_IoT: The role of heterogeneity and the need for standardization of features and attack types in IoT network intrusion data sets,'' \emph{IEEE Internet Things J.}, vol. 9, no. 1, pp. 485--496, 2022. \href{https://doi.org/10.1109/JIOT.2021.3085194}{[Online]}

\bibitem{teixeira2018scada}
M. A. Teixeira, T. Salman, M. Zolanvari, R. Jain, N. Meskin, and M. Samaka, ``SCADA system testbed for cybersecurity research using machine learning approach,'' \emph{Future Internet}, vol. 10, no. 8, Art. no. 76, 2018. \href{https://doi.org/10.3390/fi10080076}{[Online]}

\bibitem{teixeira2018wustl}
M. A. Teixeira, M. Zolanvari, and R. Jain, ``WUSTL-IIOT-2018 dataset for ICS (SCADA) cybersecurity research,'' IEEE DataPort, 2018. \href{https://doi.org/10.21227/kzgp-7t84}{[Online]}

\end{thebibliography}
\end{document}
